\begin{proof}[Proof of \cref{obj:thm:uniform-private-upper}]
\begin{enumerate}
\item
We first establish privacy of the local releases. The public scales satisfy
\[
\delta=\frac{2h}{k}>0,
\qquad
\mathcal A=nh\min\{1,n\delta\}>0,
\qquad
d_0=\frac{3\mathcal A}{64}>0.
\]
The finite cell statistics are measurable, and the positive denominator floor makes the ratio and interval maps everywhere defined and measurable. For
\[
u=(u_N,u_D)\in\mathbb R^2,
\]
the conditional density of the joint release is
\[
p_D(u)
=
\left(\frac{\epsilon}{24}\right)^2
\exp\left\{
-\frac{\epsilon}{12}
\bigl(|u_N-S_N(D)|+|u_D-S_D(D)|\bigr)
\right\}.
\]
It integrates to one for every dataset. For replacement-adjacent datasets, \cref{obj:lem:all-count-stability} supplies
\[
|S_N(D)-S_N(D')|+|S_D(D)-S_D(D')|\le12.
\]
The triangle inequality therefore gives
\[
p_D(u)\le e^\epsilon p_{D'}(u)
\qquad\text{for every }u\in\mathbb R^2.
\]
Integration proves the required privacy inequality for every Borel output event. Measurable postprocessing preserves both the Markov property and this inequality, proving the assertions for the scalar and interval outputs in \cref{obj:def:private-ratio}.
% lean: CausalSmith.Stat.PrivateCateRoughdesign.privateRatioRelease_private; CausalSmith.Stat.PrivateCateRoughdesign.Thk_private; CausalSmith.Stat.PrivateCateRoughdesign.Ihk_private

\item
The public tuning also preserves these properties. When \(r(n,\epsilon)\ge1/8\), the outputs in \cref{obj:def:public-tuning} are constant probability kernels and satisfy the privacy inequality.

On the branch \(r(n,\epsilon)<1/8\), define
\[
q_{\mathrm{dyad}}
=
\left\lceil\log_2\frac1{r(n,\epsilon)}\right\rceil,
\qquad
h_* =2^{-q_{\mathrm{dyad}}},
\qquad
k_* =2\cdot2^{4q_{\mathrm{dyad}}},
\qquad
\delta_*=\frac{2h_*}{k_*}.
\]
Here \(r(n,\epsilon)>0\), and \(q_{\mathrm{dyad}}\) is a nonnegative integer. The ceiling inequalities imply
\[
\frac1{r(n,\epsilon)}
\le2^{q_{\mathrm{dyad}}}
\le\frac2{r(n,\epsilon)},
\]
hence
\[
0<\frac{r(n,\epsilon)}2
\le h_*
\le r(n,\epsilon)<\frac18.
\]
Moreover,
\[
k_*=2h_*^{-4}\in\mathbb N,
\qquad
k_*\ge2,
\qquad
\delta_*=h_*^5.
\]
Thus the active branch uses admissible local parameters. Applying the local privacy conclusions proves the Markov and privacy assertions for both publicly tuned outputs.
% lean: CausalSmith.Stat.PrivateCateRoughdesign.tunedH_bounds; CausalSmith.Stat.PrivateCateRoughdesign.public_tuning_parameters; CausalSmith.Stat.PrivateCateRoughdesign.publicTunedRelease_private; CausalSmith.Stat.PrivateCateRoughdesign.publicTunedInterval_private

\item
It remains to establish the statistical guarantees. We reduce them to the following uniform squared-error claim: for every
\[
k'\in\mathbb N,\quad k'\ge2,
\qquad
h'\in(0,1/4],
\qquad
\delta'=\frac{2h'}{k'},
\qquad
P\in\mathcal P,
\]
one has
\[
E_{(P^{\mathrm{obs}})^{\otimes n},\widehat T_{h',k'}}
\bigl[(\widehat T_{h',k'}-\theta(P))^2\bigr]
\le V(h',\delta').
\]
We first derive the remaining conclusions from this claim, and then prove it.
% lean: CausalSmith.Stat.PrivateCateRoughdesign.uniform_private_upper

\item
For the original local parameters, define the interval radius
\[
\rho=\sqrt{10V(h,\delta)}>0.
\]
By \cref{obj:lem:point-version}, the target belongs to \([-1,1]\). Consequently, the endpoint formulas in \cref{obj:def:private-ratio} give
\[
\{\theta(P)\notin\widehat I_{h,k}\}
\subseteq
\{|\widehat T_{h,k}-\theta(P)|>\rho\}
\subseteq
\{(\widehat T_{h,k}-\theta(P))^2\ge10V(h,\delta)\}.
\]
Markov's inequality and the squared-error claim yield
\[
\Pr\{\theta(P)\notin\widehat I_{h,k}\}
\le
\frac{
E[(\widehat T_{h,k}-\theta(P))^2]
}{10V(h,\delta)}
\le\frac1{10}.
\]
All probabilities and expectations here integrate both the sample and release randomization. This proves the local coverage assertion, including targets at the endpoints of \([-1,1]\).

The interval contains its clipped center. Its endpoints lie in \([-1,1]\) and within distance \(\rho\) of that center, so, on every release,
\[
0\le\operatorname{Leb}(\widehat I_{h,k})
\le2,
\qquad
\operatorname{Leb}(\widehat I_{h,k})\le2\rho.
\]
Integrating gives
\[
E\operatorname{Leb}(\widehat I_{h,k})
\le\min\{2,2\sqrt{10V(h,\delta)}\}.
\]
% lean: CausalSmith.Stat.PrivateCateRoughdesign.Ihk_coverage_of_squared_error_le; CausalSmith.Stat.PrivateCateRoughdesign.Ihk_expectedLength_le

\item
We next prove the uniform absolute-error bound for the public estimator. If \(r(n,\epsilon)\ge1/8\), \cref{obj:lem:point-version} gives
\[
E|\widehat T^*-\theta(P)|
=|\theta(P)|
\le1
\le2^{18}r(n,\epsilon).
\]

For the active branch, the cap in the benchmark definition is inactive. Raising its three constituent lower bounds to their respective positive integer powers gives
\[
1\le nr(n,\epsilon)^4,
\qquad
1\le n^2\epsilon r(n,\epsilon)^7,
\qquad
1\le n\epsilon r(n,\epsilon)^2.
\]
In particular, since \(0<r(n,\epsilon)<1\),
\[
nr(n,\epsilon)^3\ge nr(n,\epsilon)^4\ge1,
\qquad
n^2r(n,\epsilon)^8
=\bigl(nr(n,\epsilon)^4\bigr)^2\ge1.
\]
Define the tuned occupancy scale by
\[
\mathcal A_*
=nh_*\min\{1,n\delta_*\}
=\min\{nh_*,n^2h_*^6\}.
\]
The dyadic bounds imply
\[
h_*\ge\frac{r(n,\epsilon)}{64},
\qquad
h_*^6\ge\frac{r(n,\epsilon)^6}{64},
\]
and therefore
\[
\mathcal A_*
\ge
\frac1{64}
\min\{nr(n,\epsilon),n^2r(n,\epsilon)^6\}.
\]
Multiplying by the appropriate positive scales yields
\[
\begin{aligned}
r(n,\epsilon)^2\mathcal A_*
&\ge
\frac1{64}
\min\{nr(n,\epsilon)^3,n^2r(n,\epsilon)^8\}
\ge\frac1{64},\\
r(n,\epsilon)\epsilon\mathcal A_*
&\ge
\frac1{64}
\min\{n\epsilon r(n,\epsilon)^2,
      n^2\epsilon r(n,\epsilon)^7\}
\ge\frac1{64}.
\end{aligned}
\]
Thus
\[
\mathcal A_*^{-1}\le64r(n,\epsilon)^2,
\qquad
(\epsilon\mathcal A_*)^{-1}\le64r(n,\epsilon),
\qquad
(\epsilon\mathcal A_*)^{-2}\le4096r(n,\epsilon)^2.
\]
Since \(\delta_*=h_*^5\) and \(h_*>0\),
\[
\delta_*^{2/5}=h_*^2\le r(n,\epsilon)^2.
\]
Substitution into the public envelope gives
\[
V(h_*,\delta_*)
\le
2^{20}(2+64+4096)r(n,\epsilon)^2
\le2^{33}r(n,\epsilon)^2
\le\bigl(2^{18}r(n,\epsilon)\bigr)^2.
\]

Define the positive absolute-error scale
\[
\beta_{\mathrm{abs}}=2^{18}r(n,\epsilon)>0.
\]
For every scalar output \(u_{\mathrm{sc}}\in\mathbb R\), expanding the nonnegative square
\((|u_{\mathrm{sc}}-\theta(P)|-\beta_{\mathrm{abs}})^2\) gives
\[
|u_{\mathrm{sc}}-\theta(P)|
\le
\frac{(u_{\mathrm{sc}}-\theta(P))^2}{2\beta_{\mathrm{abs}}}
+\frac{\beta_{\mathrm{abs}}}{2}.
\]
The uniform squared-error claim applies at \((h_*,k_*)\). Integrating the last inequality therefore proves
\[
E|\widehat T^*-\theta(P)|
\le
\frac{V(h_*,\delta_*)}{2\beta_{\mathrm{abs}}}
+\frac{\beta_{\mathrm{abs}}}{2}
\le\beta_{\mathrm{abs}}.
\]
Taking the supremum over \(P\in\mathcal P\) proves the stated risk bound.
% lean: CausalSmith.Stat.PrivateCateRoughdesign.active_rate_constraints; CausalSmith.Stat.PrivateCateRoughdesign.tuned_info_lower; CausalSmith.Stat.PrivateCateRoughdesign.tuned_scaled_info_lower; CausalSmith.Stat.PrivateCateRoughdesign.tuned_inverse_info_upper; CausalSmith.Stat.PrivateCateRoughdesign.tuned_cell_roughness; CausalSmith.Stat.PrivateCateRoughdesign.tuned_Vbound_le; CausalSmith.Stat.PrivateCateRoughdesign.absolute_moment_le_of_squared_moment; CausalSmith.Stat.PrivateCateRoughdesign.publicTunedRelease_worstRisk_of_squared_error_le

\item
For the public interval, the fallback branch has deterministic length two, and
\[
E\operatorname{Leb}(\widehat I^*)=2
\le2^{22}r(n,\epsilon)
\qquad\text{when }r(n,\epsilon)\ge1/8.
\]
On the active branch, the envelope estimate just proved gives
\[
10V(h_*,\delta_*)
\le10\cdot2^{33}r(n,\epsilon)^2
\le\bigl(2^{21}r(n,\epsilon)\bigr)^2.
\]
The right-hand square root is nonnegative, so
\[
\sqrt{10V(h_*,\delta_*)}\le2^{21}r(n,\epsilon).
\]
The local expected-length bound consequently yields
\[
E\operatorname{Leb}(\widehat I^*)
\le2\sqrt{10V(h_*,\delta_*)}
\le2^{22}r(n,\epsilon).
\]
Taking the supremum over the model proves the public expected-length assertion.
% lean: CausalSmith.Stat.PrivateCateRoughdesign.tuned_interval_radius_le; CausalSmith.Stat.PrivateCateRoughdesign.publicTunedInterval_expectedLength_le; CausalSmith.Stat.PrivateCateRoughdesign.publicTunedInterval_worstLength_le

\item
Public coverage follows through the same two tuning branches. On the fallback branch, \cref{obj:lem:point-version} gives
\[
\Pr\{\theta(P)\in[-1,1]\}=1.
\]
On the active branch, the admissibility of \(h_*,k_*\) allows application of the uniform squared-error claim:
\[
E[(\widehat T_{h_*,k_*}-\theta(P))^2]
\le V(h_*,\delta_*).
\]
The local coverage argument then gives
\[
\Pr\{\theta(P)\in\widehat I^*\}
=
\Pr\{\theta(P)\in\widehat I_{h_*,k_*}\}
\ge\frac9{10}.
\]
% lean: CausalSmith.Stat.PrivateCateRoughdesign.publicTunedInterval_coverage_of_squared_error_le

\item
We now prove the uniform squared-error claim. Fix arbitrary admissible local parameters \(h,k\) and \(P\in\mathcal P\). For measurable integrands, write
\[
\mathbb E_P[F(D,\widetilde S)]
:=
\int_{\mathcal O^n}\int_{\mathbb R^2}
F(D,u)\,\widetilde S(D,du)\,
(P^{\mathrm{obs}})^{\otimes n}(dD).
\]
Thus this expectation includes the independent Laplace randomization. Define the population means, bias envelope, and comparison ratio by
\[
\overline N=\mathbb E_P[S_N(D)],
\qquad
\overline D=\mathbb E_P[S_D(D)],
\qquad
b=3h+24\delta^{1/5},
\qquad
\theta^\circ=\frac{\overline N}{\overline D}.
\]
By \cref{obj:lem:occupancy-cancellation},
\[
\overline N=\sum_{j=1}^k w_j\nu_j,
\qquad
\overline D=\sum_{j=1}^k w_jd_j,
\qquad
\overline D\ge d_0>0,
\qquad
|\theta^\circ-\theta(P)|\le b.
\]

For binary values
\(a_{\mathrm{obs}},a'_{\mathrm{obs}},
y_{\mathrm{obs}},y'_{\mathrm{obs}}\in\{0,1\}\),
\[
\left|
(a_{\mathrm{obs}}-a'_{\mathrm{obs}})
(y_{\mathrm{obs}}-y'_{\mathrm{obs}})
\right|
\le
(a_{\mathrm{obs}}-a'_{\mathrm{obs}})^2.
\]
Taking expectations under two independent observations conditional on cell \(j\) proves
\[
|\nu_j|\le d_j.
\]
The cell conditioning is well defined by the positive cell masses under the density bounds. The occupancy weights in \cref{obj:lem:occupancy-cancellation} are nonnegative, hence
\[
|\overline N|
\le\sum_{j=1}^k w_j|\nu_j|
\le\sum_{j=1}^k w_jd_j
=\overline D,
\qquad
|\theta^\circ|\le1.
\]
Also, \cref{obj:lem:point-version} supplies
\[
\theta(P)\in[-1,1].
\]
% lean: CausalSmith.Stat.PrivateCateRoughdesign.conditional_cell_numerator_abs_le_denominator; CausalSmith.Stat.PrivateCateRoughdesign.population_numerator_abs_le_denominator; CausalSmith.Stat.PrivateCateRoughdesign.population_ratio_abs_le_one

\item
For every possible joint output \(u=(u_N,u_D)\in\mathbb R^2\), define
\[
B_D(u)=\max\{d_0,u_D\},
\qquad
T_{\mathrm{loc}}(u)
=
\operatorname{clip}_{[-1,1]}
\left(\frac{u_N}{B_D(u)}\right).
\]
These maps are defined on all of \(\mathbb R^2\), with \(B_D(u)\ge d_0>0\). Since \(\overline D\ge d_0\), the denominator floor fixes \(\overline D\) and satisfies
\[
|B_D(u)-\overline D|\le|u_D-\overline D|.
\]
The exact perturbation identity is
\[
\frac{u_N}{B_D(u)}-\theta^\circ
=
\frac{
(u_N-\overline N)+
\theta^\circ(\overline D-B_D(u))
}{B_D(u)}.
\]
Using \(|\theta^\circ|\le1\) gives
\[
\left|\frac{u_N}{B_D(u)}-\theta^\circ\right|
\le
\frac{|u_N-\overline N|+|u_D-\overline D|}{d_0}.
\]
Clipping cannot increase distance to \(\theta(P)\in[-1,1]\). Combining this fact with the population bias bound yields
\[
|T_{\mathrm{loc}}(u)-\theta(P)|
\le
b+\frac{|u_N-\overline N|+|u_D-\overline D|}{d_0}.
\]
Two applications of the quadratic sum inequality now give
\[
\begin{aligned}
(T_{\mathrm{loc}}(u)-\theta(P))^2
&\le
2b^2+
\frac{2}{d_0^2}
\bigl(|u_N-\overline N|+|u_D-\overline D|\bigr)^2\\
&\le
2b^2+\frac4{d_0^2}
\left\{
(u_N-\overline N)^2+(u_D-\overline D)^2
\right\}.
\end{aligned}
\]
% lean: CausalSmith.Stat.PrivateCateRoughdesign.floored_ratio_perturbation; CausalSmith.Stat.PrivateCateRoughdesign.ratioMap_error_le; CausalSmith.Stat.PrivateCateRoughdesign.ratioMap_squared_error_le; CausalSmith.Stat.PrivateCateRoughdesign.ratioMap_model_squared_error_le

\item
For the variance calculation, invoke the replacement-copy Efron--Stein inequality \citep[Theorem~3.1, replacement-copy form, p.~54]{BoucheronLugosiMassart2013}. Explicitly, for the independent observed records and their independent replacement copies, define
\[
D=(O_1,\ldots,O_n),
\qquad
D^{(i)}=(O_1,\ldots,O_{i-1},O_i',O_{i+1},\ldots,O_n),
\]
where \(O_1,\ldots,O_n,O_i'\) have law \(P^{\mathrm{obs}}\), and \(O_i'\) is independent of \(D\). The published inequality gives
\[
\operatorname{Var}(\mathsf S(D))
\le
\frac12\sum_{i=1}^n
E\bigl[(\mathsf S(D)-\mathsf S(D^{(i)}))^2\bigr],
\qquad
\mathsf S\in\{S_N,S_D\}.
\]
The statistics are measurable and square integrable. With this inequality, \cref{obj:lem:all-count-stability} supplies the bound
\[
\max\{\operatorname{Var}(S_N(D)),\operatorname{Var}(S_D(D))\}
\le18W(P),
\qquad
W(P)=\sum_{j=1}^k w_j.
\]

Each independent Laplace noise has
\[
E[\xi_N]=E[\xi_D]=0,
\qquad
E[\xi_N^2]=E[\xi_D^2]
=2\left(\frac{12}{\epsilon}\right)^2
=\frac{288}{\epsilon^2}.
\]
Expanding the centered squares, independence and the zero noise means eliminate the cross terms. Thus the exact released-moment identity is
\[
\begin{aligned}
&\mathbb E_P\left[
(\widetilde S_N-\overline N)^2+
(\widetilde S_D-\overline D)^2
\right]\\
&\qquad=
\operatorname{Var}(S_N(D))
+\operatorname{Var}(S_D(D))
+\frac{576}{\epsilon^2}
\le36W(P)+\frac{576}{\epsilon^2}.
\end{aligned}
\]
Integrating the pointwise ratio bound consequently gives
\[
\begin{aligned}
\mathbb E_P[(\widehat T_{h,k}-\theta(P))^2]
&\le
2b^2+\frac4{d_0^2}
\mathbb E_P\left[
(\widetilde S_N-\overline N)^2+
(\widetilde S_D-\overline D)^2
\right]\\
&\le
2b^2+\frac4{d_0^2}
\left(36W(P)+\frac{576}{\epsilon^2}\right)\\
&=
2b^2+\frac{144W(P)}{d_0^2}
+\frac{2304}{\epsilon^2d_0^2}.
\end{aligned}
\]
% lean: CausalSmith.Stat.PrivateCateRoughdesign.privateRatioRelease_centered_moments; CausalSmith.Stat.PrivateCateRoughdesign.Thk_model_squared_error_le

\item
Finally, the quadratic sum inequality and positivity of \(\delta\) imply
\[
(\delta^{1/5})^2=\delta^{2/5},
\qquad
2b^2
=2(3h+24\delta^{1/5})^2
\le36h^2+2304\delta^{2/5}.
\]
By \cref{obj:lem:occupancy-cancellation},
\[
W(P)\le\frac92\mathcal A.
\]
Substituting \(d_0=3\mathcal A/64\) gives
\[
\frac{144W(P)}{d_0^2}
\le
\frac{144(9\mathcal A/2)}{(3\mathcal A/64)^2}
=294912\mathcal A^{-1},
\]
and
\[
\frac{2304}{\epsilon^2d_0^2}
=
\frac{2304}{\epsilon^2(3\mathcal A/64)^2}
=1048576(\epsilon\mathcal A)^{-2}.
\]
All four terms in the public envelope are nonnegative. Since
\(36,2304,294912\le2^{20}\) and \(1048576=2^{20}\), we obtain
\[
\begin{aligned}
\mathbb E_P[(\widehat T_{h,k}-\theta(P))^2]
&\le
36h^2+2304\delta^{2/5}
+294912\mathcal A^{-1}
+1048576(\epsilon\mathcal A)^{-2}\\
&\le
2^{20}\left\{
h^2+\delta^{2/5}+\mathcal A^{-1}
+(\epsilon\mathcal A)^{-2}
\right\}
=V(h,\delta).
\end{aligned}
\]
The local parameters and model law were arbitrary, proving the uniform squared-error claim and completing all the preceding deductions.
% lean: CausalSmith.Stat.PrivateCateRoughdesign.bias_squared_le; CausalSmith.Stat.PrivateCateRoughdesign.variance_floor_envelope_le; CausalSmith.Stat.PrivateCateRoughdesign.model_variance_envelope_le; CausalSmith.Stat.PrivateCateRoughdesign.Thk_model_squared_error_le; CausalSmith.Stat.PrivateCateRoughdesign.uniform_private_upper
\end{enumerate}
\end{proof}
